\documentclass[11pt,reqno]{amsart}

%==============================================================
% The Even Pseudo-Process Bridge to the Riemann Hypothesis:
% the Endpoint Dichotomy and the Moment-Oscillation Theorem
%==============================================================

\usepackage[utf8]{inputenc}
\usepackage[T1]{fontenc}
\usepackage{lmodern}
\usepackage[margin=1.15in]{geometry}
\usepackage{amsmath,amssymb,amsthm,mathtools}
\usepackage{booktabs}
\usepackage{microtype}
\usepackage{hyperref}
\usepackage{cleveref}
\emergencystretch=3em

\theoremstyle{plain}
\newtheorem{theorem}{Theorem}[section]
\newtheorem{proposition}[theorem]{Proposition}
\newtheorem{lemma}[theorem]{Lemma}
\newtheorem{corollary}[theorem]{Corollary}
\theoremstyle{definition}
\newtheorem{definition}[theorem]{Definition}
\newtheorem{hypothesis}[theorem]{Hypothesis}
\newtheorem{problem}[theorem]{Open problem}
\theoremstyle{remark}
\newtheorem{remark}[theorem]{Remark}

\newcommand{\R}{\mathbb R}
\newcommand{\C}{\mathbb C}
\newcommand{\N}{\mathbb N}
\newcommand{\Z}{\mathbb Z}
\newcommand{\E}{\mathbb E}
\newcommand{\LP}{\mathcal{LP}}
\newcommand{\T}{\mathcal T}
\newcommand{\Mcal}{\mathcal M}
\DeclareMathOperator{\sgn}{sgn}
\DeclareMathOperator{\Log}{Log}

\title[Failure of the even-cone bridge to RH]{Moment Oscillation and the
Failure of a One-Sided Even Pseudo-Process Cone Bridge\\ to the Riemann
Hypothesis}
\author{Ramiro Fontes}
\address{Quijotic Research}
\email{ramirofontes@gmail.com}
\urladdr{www.quijotic.com}
\date{\today}

\begin{document}

\begin{abstract}
The de Bruijn--Newman flow $H_{t}(x)=\int_0^\infty
e^{tu^2}\Phi(u)\cos(xu)\,du$ embeds in the symbol family
$H_{t,\alpha}=\int_0^\infty e^{tu^\alpha}\Phi(u)\cos(xu)\,du$, and the
Riemann hypothesis is $\Lambda\le0$ for its real-rooting threshold. The
companion paper~\cite{Fsymbol} closed the first-order route along the
\emph{logarithmic} tangent $u^2\log u=\partial_\alpha u^\alpha|_{\alpha=2}$:
the self-dual extremality hypothesis fails because the tilt numerator is
strictly positive. This paper examines the other first-order route a reader
naturally tries, the one carried by the \emph{even integer} exponents
$\alpha=2m$, $m=2,3,4,\dots$, the backward higher-order heat
pseudo-processes of Krylov--Hochberg. Perturbing $u^2\rightsquigarrow
u^2+\epsilon u^{2m}$ at a nondegenerate positive-$\Lambda$ collision
$(x_*,\Lambda)$ moves the collision time at the explicit rate
$T_m'(0)=-M_m/M_1$, with $M_m:=\int_0^\infty
u^{2m}e^{\Lambda u^2}\Phi(u)\cos(x_*u)\,du$. A one-sided ``cone'' bridge to
RH would require the response moments $M_m$ to be eventually of one sign;
the bridge breaks because they are not. We prove this through a single
structural mechanism, the \emph{endpoint dichotomy}. Shifting the cosine
transform to the contour height $h\in(0,\pi/8)$, the forced sign of a
response is carried entirely by the imaginary part of the symbol on the
ray $u=iy$: for the logarithmic weight that imaginary part is $-\tfrac\pi2
y^2$, producing the positive Laplace transform of~\cite{Fsymbol}; for every
pure even power $u^{2m}$ the symbol $(iy)^{2m}=(-1)^my^{2m}$ is \emph{real},
the vertical segment contributes nothing, and $M_m$ is left as the bare
oscillatory horizontal ray. The same imaginary-part density
$y^\alpha\cos\tfrac{\pi(\alpha+1)}2$ governs all three phenomena of the
program at once: it is the infrared-catastrophe coefficient of the odd
exponents, its $\alpha$-derivative at $2$ is the tilt positivity, and its
vanishing on the even integers is precisely the present oscillation. We
then prove the \emph{moment-oscillation theorem}: for $x_*\neq0$ the
sequence $M_m$ obeys $M_m=A_m\cos(x_*u_m+\theta_m)(1+o(1))$ with $A_m>0$,
$\theta_m\to0$ and the arithmetic saddle $u_m=\tfrac14\log m-\tfrac14\log\log
m+O(1)\to\infty$, so $M_m$ changes sign infinitely often; in particular
$\{M_m\}$ is neither eventually one-signed nor a Stieltjes moment sequence,
which closes the one-sided even-cone bridge in both its eventual-sign and
positive-moment forms --- not every conceivable first-order mechanism, but
the natural cone a reader is led to try. Numerics on the Riemann kernel
confirm the oscillation: at the surrogate collision sites
$x_*=2\gamma_1\approx28.27$ and $x_*=2\gamma_3\approx50.02$ the velocity
sign $\sgn(-M_m/M_1)$ reverses four and six times across $m\le30$, faster at
the larger site exactly as the saddle requires, and stably under the flow
$t$. The level-$p$ Clark--Ocone calculus of~\cite{Fpseudo} supplies the
rigorous language for the response-measure pairing and, more pointedly, the
reason the route cannot succeed: for the orders $p=2m\ge4$ that would have to
supply the cone there is no positive path-space law, no orthogonal chaos,
and no Laguerre--P\'olya preservation; the pseudo-process's own moments are
sign-regular, $\E[X_T^{pn}]=\tfrac{(pn)!}{n!}c_p^nT^n$, but the
\emph{arithmetic} response moments $M_m$ carry the zero's phase
$\cos(x_*u)$ and oscillate, and these are moments of different measures
between which the calculus does not transfer. The first-order even bridge
therefore closes; what survives is not a pointwise sign but a global
analytic question on the function $\Mcal(\beta)=\int_0^\infty
u^{2\beta}e^{\Lambda u^2}\Phi\cos(x_*u)\,du$, and the one even exponent that
keeps genuine positivity, the quartic $\alpha=4$ iterated-Brownian flow,
whose global real-rootedness is the surviving frontier.
\end{abstract}

\maketitle

%==============================================================
\section{Introduction}\label{sec:intro}
%==============================================================

\subsection*{The flow, the constant, and the two first-order routes}
The Riemann $\xi$-function has the Fourier representation
\begin{equation}\label{eq:xi-rep}
  \tfrac18\,\xi\!\left(\tfrac12+\tfrac{ix}{2}\right)
  =\int_0^\infty \Phi(u)\cos(xu)\,du,
  \quad
  \Phi(u)=\sum_{n\ge1}\bigl(2\pi^2n^4e^{9u}-3\pi n^2e^{5u}\bigr)
  e^{-\pi n^2e^{4u}},
\end{equation}
with $\Phi$ positive, even, and decaying like $\exp(-\pi e^{4u})$; the
evenness is the theta modularity~\cite{Titchmarsh,Polya27}. The de
Bruijn--Newman flow~\cite{deBruijn,Newman}
\begin{equation}\label{eq:dbn}
  H_t(x):=\int_0^\infty e^{tu^2}\Phi(u)\cos(xu)\,du
\end{equation}
has only real zeros exactly for $t\ge\Lambda$, the Riemann hypothesis is
$\Lambda\le0$, and $\Lambda\ge0$ is known~\cite{RodgersTao}, so RH is the
equality $\Lambda=0$: the flow is critical now, and a positive $\Lambda$
would be attained at a real collision of two zeros confined to a compact
set (effective Ki--Kim--Lee, \cite{Polymath15}, and the companion's
\cite{Fsymbol} compact-criticality theorem).

The companion paper~\cite{Fsymbol} deforms the \emph{generator} rather than
the time, through
\begin{equation}\label{eq:family}
  H_{t,\alpha}(x):=\int_0^\infty e^{tu^\alpha}\Phi(u)\cos(xu)\,du,
  \qquad (t,\alpha)\in\R\times(0,\infty),
\end{equation}
so $H_{t,2}=H_t$. At a real collision $x_*$ of multiplicity exactly two and
time $\Lambda$, the collision time $T_*(\alpha)$ tracked through the family
is real-analytic near $\alpha=2$ and moves at the rate $T_*'(2)=-\Lambda
B_*$, where
\begin{equation}\label{eq:Bstar}
  B_*=\frac{N(x_*,\Lambda)}{D(x_*,\Lambda)},\quad
  D(x,t)=\!\int_0^\infty\! u^2e^{tu^2}\Phi\cos(xu)\,du,\ \
  N(x,t)=\!\int_0^\infty\! u^2\log u\,e^{tu^2}\Phi\cos(xu)\,du,
\end{equation}
and $D=-\partial_x^2H$. The hoped-for ``self-dual extremality'' (SD) made
$\alpha=2$ a local minimum of $T_*$, forcing $T_*'(2)=0$ and hence (with
$\Lambda>0$) a contradiction --- that was the route to RH. The
companion~\cite{Fsymbol} \emph{refutes} (SD): the tilt numerator
$N(x,t)>0$ for all $x\ge27$, $t\in[0,\tfrac12]$, by a contour shift, so
$B_*\neq0$ and the first-order route along the \emph{logarithmic} tangent
$u^2\log u=\partial_\alpha u^\alpha|_{\alpha=2}$ closes negatively.

That tangent is one direction in the family. A reader who has seen it close
asks the obvious next question. The even integers $\alpha\in\{4,6,8,\dots\}$
are the only \emph{other} exponents in \eqref{eq:family} that do not
instantly destroy real-rootedness (the infrared catastrophe
of~\cite{Fsymbol} makes $\Lambda(\alpha)=+\infty$ for every
$\alpha\notin2\N$); they are the backward higher-order heat
\emph{pseudo-processes} of Krylov and Hochberg~\cite{Krylov,Hochberg},
whose representation calculus is the companion~\cite{Fpseudo}. Does the
\emph{even pseudo-process hierarchy} open a route that the single
logarithmic tangent could not?

\subsection*{The even pseudo-process bridge}
Perturb the heat generator by a higher even power,
\begin{equation}\label{eq:Gm}
  G^{(m)}_{s,\epsilon}(x):=\int_0^\infty
  e^{su^2+\epsilon u^{2m}}\,\Phi(u)\cos(xu)\,du,\qquad m\ge1,
\end{equation}
and track the collision. The first-order motion is the response moment
ratio
\begin{equation}\label{eq:Tm}
  T_m'(0)=-\frac{M_m}{M_1},\qquad
  M_m:=\int_0^\infty u^{2m}e^{\Lambda u^2}\Phi(u)\cos(x_*u)\,du,
\end{equation}
with $M_1=D(x_*,\Lambda)=-H_\Lambda''(x_*)\neq0$
(\cref{thm:velocity}). The bridge a reader hopes for is a \emph{cone}: that
every higher even direction pushes the collision the same way, i.e.\ that
$\sgn(T_m'(0))$ --- equivalently $\sgn(M_m)$ --- is eventually constant in
$m$. A fixed sign would say the bad collision is variationally one-sided
against the whole even cone, which, combined with the oscillatory growth of
the moment order, is the contradiction sought. The entire question is the
sign behavior of the sequence $\{M_m\}$.

\subsection*{Three results, and the verdict}
The verdict is that the bridge closes, and the reason is a single mechanism.
\begin{enumerate}
\item \textbf{The collision velocity} (\cref{sec:velocity}). At a
nondegenerate positive-$\Lambda$ collision the even direction $u^{2m}$ moves
the collision time at the rate \eqref{eq:Tm}; the derivation is the
implicit function theorem of~\cite{Fsymbol} with $\epsilon$ in place of
$\alpha$.
\item \textbf{The endpoint dichotomy} (\cref{sec:dichotomy}). Shift the
cosine transform to the contour $[0,ih]\cup[ih,ih+\infty)$, $h\in(0,\pi/8)$.
The forced part of a response is the real part of the vertical segment, and
that real part is the \emph{imaginary part of the symbol on the imaginary
axis}: for $u^\alpha$ it is the density $y^\alpha\cos\tfrac{\pi(\alpha+1)}2$.
This one density runs the whole program: it is the
infrared-catastrophe coefficient of~\cite{Fsymbol} for $\alpha\notin2\N$,
its $\alpha$-derivative at $\alpha=2$ is the tilt positivity
$+\tfrac\pi2y^2$ of~\cite{Fsymbol}, and it \emph{vanishes identically on the
even integers}. Hence the even moments $M_m$ have no forced segment term at
all (\cref{thm:dichotomy}); they are carried entirely by the oscillatory
ray.
\item \textbf{The moment-oscillation theorem} (\cref{sec:oscillation}). For
$x_*\neq0$,
\[
   M_m=A_m\cos(x_*u_m+\theta_m)\,(1+o(1)),\quad A_m>0,\ \theta_m\to0,\
   u_m=\tfrac14\log m-\tfrac14\log\log m+O(1)\to\infty,
\]
so $M_m$ changes sign infinitely often (\cref{thm:oscillation}). The
sequence is therefore neither eventually one-signed nor a Stieltjes moment
sequence: the one-sided even-cone bridge fails in both its eventual-sign and
positive-moment forms (\cref{cor:nogo}).
\end{enumerate}
\Cref{sec:closure} assembles the closure and unifies it with the (SD)
refutation through the response measure $\mu_*$ and the function
$R(\beta)=\int u^{2\beta}d\mu_*$: the value $R(1)=1$ and the slope
$R'(1)=2B_*>0$ are \emph{forced} (the trivial and logarithmic directions,
both endpoint-singular), while the integer samples $R(m)$ \emph{oscillate}
(the even pseudo-process directions, endpoint-smooth). The discriminant is
endpoint regularity, nothing else. \Cref{sec:probabilistic} reads the
result through the level-$p$ Clark--Ocone calculus~\cite{Fpseudo}, which
supplies both the representation language and the structural reason the cone
cannot exist. \Cref{sec:numerics} reports the confirming computation.
\Cref{sec:frontier} states what survives: not a pointwise sign but a global
analytic question, and the quartic flow.

\begin{center}
\fbox{\begin{minipage}{0.93\linewidth}\small
\textbf{What this paper settles.} The first-order route to RH along the even
pseudo-process hierarchy $\alpha=2m$ closes, because the response moments
$M_m$ oscillate in sign. The oscillation is not arithmetic bad luck: it is
forced by the same contour mechanism that, on the singular directions,
forces a definite sign. The level-$p$ Clark--Ocone
calculus~\cite{Fpseudo} confirms the obstruction --- the orders $p=2m\ge4$
have no positivity to donate --- rather than removing it. Together
with~\cite{Fsymbol}, both first-order deformation routes (logarithmic and
even-integer) are now resolved negatively: one because its invariant is too
positive to vanish, the other because its invariant has no fixed sign.
\end{minipage}}
\end{center}

%==============================================================
\section{The collision velocity along even pseudo-process directions}
\label{sec:velocity}
%==============================================================

We work at a base collision furnished by the unconditional compact
criticality of the companion paper. Recall from~\cite{Fsymbol} that if
$\Lambda>0$ then $H_\Lambda$ has a real zero $x_*$ of multiplicity at least
two with $|x_*|>27.2$, confined to a compact set; nondegeneracy (ND) is the
statement that the multiplicity is exactly two, i.e.\ $H_\Lambda''(x_*)\neq0$.
All integrals below converge absolutely and may be differentiated under the
integral sign because $\Phi$ decays doubly exponentially while
$e^{su^2+\epsilon u^{2m}}$ grows only like $\exp(O(u^{2m}))$; the data
$(s,\epsilon,x)\mapsto G^{(m)}_{s,\epsilon}(x)$ of \eqref{eq:Gm} are jointly
real-analytic (the well-posedness lemma of~\cite{Fsymbol} applies verbatim,
$u^{2m}$ being a smooth even symbol).

\begin{theorem}[Even-direction collision velocity]\label{thm:velocity}
Suppose $\Lambda>0$ and $H_\Lambda$ has a real zero $x_*$ of multiplicity
exactly two. Fix an integer $m\ge1$. There are real-analytic functions
$\epsilon\mapsto(x_m(\epsilon),T_m(\epsilon))$ near $\epsilon=0$, with
$(x_m(0),T_m(0))=(x_*,\Lambda)$, such that $G^{(m)}_{T_m(\epsilon),\epsilon}$
has a real zero of multiplicity two at $x_m(\epsilon)$, and
\begin{equation}\label{eq:velocity}
  T_m'(0)=-\frac{M_m}{M_1},
  \qquad
  M_m=\int_0^\infty u^{2m}e^{\Lambda u^2}\Phi(u)\cos(x_*u)\,du,
\end{equation}
where $M_1=D(x_*,\Lambda)=-H_\Lambda''(x_*)\neq0$.
\end{theorem}

\begin{proof}
Apply the implicit function theorem to the map
$(x,s)\mapsto\bigl(G^{(m)}_{s,\epsilon}(x),\partial_xG^{(m)}_{s,\epsilon}(x)\bigr)$
at $(x_*,\Lambda,0)$. At the base point $G^{(m)}_{\Lambda,0}=H_\Lambda$, so
$G=\partial_xG=0$ and $\partial_x^2G=H_\Lambda''(x_*)\neq0$. The
$(x,s)$-Jacobian is
\[
  \det\begin{pmatrix}\partial_xG & \partial_sG\\
  \partial_x^2G & \partial_x\partial_sG\end{pmatrix}
  =\det\begin{pmatrix}0 & \partial_sG\\ \partial_x^2G & *\end{pmatrix}
  =-\,\partial_sG\,\partial_x^2G .
\]
Now $\partial_sG^{(m)}_{\Lambda,0}(x_*)=\int u^2e^{\Lambda u^2}\Phi\cos(x_*u)\,du
=D(x_*,\Lambda)=-H_\Lambda''(x_*)$ by the heat equation
$\partial_sH=-\partial_x^2H$ at $\alpha=2$, so the Jacobian equals
$(H_\Lambda''(x_*))^2\neq0$ and the collision manifold
$(x_m(\epsilon),T_m(\epsilon))$ exists and is real-analytic. Differentiating
$G^{(m)}_{T_m(\epsilon),\epsilon}(x_m(\epsilon))=0$ in $\epsilon$ and using
$\partial_xG=0$ at the collision to kill the $x_m'(\epsilon)$ term gives
$\partial_sG\cdot T_m'(0)+\partial_\epsilon G=0$. Since
$\partial_\epsilon G^{(m)}_{\Lambda,0}(x_*)=\int u^{2m}e^{\Lambda u^2}\Phi
\cos(x_*u)\,du=M_m$ and $\partial_sG=M_1$, this is \eqref{eq:velocity}.
\end{proof}

\begin{remark}[The bridge in one line]\label{rem:bridge}
By \eqref{eq:velocity} the sign of the collision velocity along $u^{2m}$ is
$-\sgn(M_m/M_1)$. A one-sided cone --- all higher even directions pushing
the same way --- is exactly the statement that $\sgn(M_m)$ is eventually
constant. The remainder of the paper computes that sign and finds it is
not.
\end{remark}

The companion's general response calculus~\cite{Fsymbol} packages all
directions at once. For a symbol direction $\sigma$ of polynomial growth,
the collision responds at the rate
$\frac{d}{d\epsilon}T_*|_\sigma=-\Lambda\int_0^\infty\sigma\,d\mu_*$ against
the single signed \emph{response measure}
\begin{equation}\label{eq:mustar}
  d\mu_*(u):=\frac{e^{\Lambda u^2}\Phi(u)\cos(x_*u)}{D(x_*,\Lambda)}\,du,
  \qquad \int_0^\infty u^2\,d\mu_*=1 .
\end{equation}
\Cref{thm:velocity} is the case $\sigma=u^{2m}$:
$\int u^{2m}d\mu_*=M_m/M_1$. The measure $\mu_*$ is signed precisely because
of the factor $\cos(x_*u)$, and that sign is the whole story.

%==============================================================
\section{The endpoint dichotomy}\label{sec:dichotomy}
%==============================================================

This section isolates the mechanism. Both the companion's tilt positivity
and the present moment oscillation are decided by what the symbol does on
the imaginary axis, and they sit at the two ends of one explicit density.

\subsection{The imaginary-axis density}
Write $g_t(u):=e^{tu^2}\Phi(u)$. By the theta modularity $\Phi$ is even and
real on $\R$ and extends holomorphically to the strip $|\Im u|<\pi/8$, where
$\Re(\pi n^2e^{4u})>0$; by Schwarz reflection $\Phi(iy)$ is real for real
$y$ with $|y|<\pi/8$, hence so is $g_t(iy)=e^{-ty^2}\Phi(iy)$. For a weight
$w$ smooth on $(0,\infty)$ consider the response transform
\begin{equation}\label{eq:Iw}
  I_w(x;t):=\int_0^\infty w(u)\,g_t(u)\cos(xu)\,du
  =\Re\int_0^\infty w(u)\,g_t(u)\,e^{ixu}\,du ,
\end{equation}
the second equality holding when $w$ is real on $(0,\infty)$. Fix
$h\in(0,\pi/8)$ and deform the ray $[0,\infty)$ to the vertical segment
$[0,ih]$ followed by the horizontal ray $[ih,ih+\infty)$; the integrand is
holomorphic between the contours and decays as $\Re u\to\infty$ because
$g_t$ super-decays and $|e^{ixu}|=e^{-x\Im u}\le1$. Thus
\begin{equation}\label{eq:split}
  I_w(x;t)=\Re\!\int_{[0,ih]}\!w\,g_t\,e^{ixu}\,du
          +\Re\!\int_{[ih,ih+\infty)}\!w\,g_t\,e^{ixu}\,du .
\end{equation}
On the segment $u=iy$, $du=i\,dy$, $e^{ixu}=e^{-xy}$, and $g_t(iy)$ is real,
so the segment integrand is $i\,w(iy)\,g_t(iy)e^{-xy}\,dy$ and
\begin{equation}\label{eq:segment}
  \Re\!\int_{[0,ih]}w\,g_t\,e^{ixu}\,du
  =-\int_0^h \bigl(\Im\,w(iy)\bigr)\,g_t(iy)\,e^{-xy}\,dy .
\end{equation}
The segment --- the only part of \eqref{eq:split} that is not exponentially
small in $x$ over the whole admissible range, the ray being
$O(e^{-hx})$ --- sees the weight only through $\Im\,w(iy)$. This is the
mechanism.

\begin{lemma}[Imaginary-axis density of the pure power]\label{lem:density}
For $w(u)=u^\alpha$ (principal branch), $\Im\,w(iy)=y^\alpha
\sin\tfrac{\pi\alpha}2$, so the segment density of \eqref{eq:segment} is
\begin{equation}\label{eq:density}
  -\Im\,(iy)^\alpha=-y^\alpha\sin\tfrac{\pi\alpha}2,
  \qquad\text{equivalently, after the factor }i:\quad
  y^\alpha\cos\tfrac{\pi(\alpha+1)}2 .
\end{equation}
In particular $\Im\,(iy)^{2m}=0$ for every integer $m$: a pure even power is
real on the imaginary axis and contributes \emph{nothing} to the segment.
For the logarithmic weight $w(u)=u^2\log u$,
$w(iy)=(iy)^2\Log(iy)=-y^2(\log y+\tfrac{i\pi}2)$, so $\Im\,w(iy)=-\tfrac\pi2
y^2$ and the segment density is $+\tfrac\pi2 y^2$.
\end{lemma}

\begin{proof}
$(iy)^\alpha=y^\alpha e^{i\pi\alpha/2}$ gives the first line; multiplying by
the contour factor $i=e^{i\pi/2}$ turns the exponent into
$\tfrac{\pi(\alpha+1)}2$, whose cosine is the displayed density and matches
the catastrophe normalization of~\cite{Fsymbol}. The logarithmic line is
the stated computation, the branch of $\Log$ taken off $(-\infty,0]$ so that
$\Log(iy)=\log y+\tfrac{i\pi}2$ for $y>0$.
\end{proof}

\subsection{One density, three phenomena}
The density $\kappa(\alpha,y):=y^\alpha\cos\tfrac{\pi(\alpha+1)}2$ of
\eqref{eq:density} is the common root of the deformation program.

\begin{enumerate}
\item[(i)] \emph{Infrared catastrophe} (\cite{Fsymbol}). For
$\alpha\notin2\N$, $\kappa(\alpha,\cdot)\not\equiv0$, and the segment
produces the sign-definite power tail
$H_{t,\alpha}(x)=t\,\Phi(0)\,\Gamma(\alpha+1)\cos\tfrac{\pi(\alpha+1)}2\,
x^{-\alpha-1}+\cdots$; the coefficient is the value
$\cos\tfrac{\pi(\alpha+1)}2$ that vanishes exactly on the even integers, so
real-rooting deformation lives only on $2\N$.
\item[(ii)] \emph{Tilt positivity} (\cite{Fsymbol}). The logarithmic tangent
is $u^2\log u=\partial_\alpha u^\alpha|_{\alpha=2}$, and
$\partial_\alpha\kappa(\alpha,y)|_{\alpha=2}
=y^2\bigl[\log y\cos\tfrac{3\pi}2-\tfrac\pi2\sin\tfrac{3\pi}2\bigr]
=+\tfrac\pi2y^2$, the positive Laplace density that proves $N(x,t)>0$ and
refutes (SD).
\item[(iii)] \emph{Moment oscillation} (this paper). The pure even powers
sit \emph{at the zeros} of $\kappa$: $\kappa(2m,\cdot)\equiv0$. They have no
segment term, hence no forced sign; their response is the horizontal ray
alone, and \cref{sec:oscillation} shows it oscillates.
\end{enumerate}

The three are one statement read at three values of $\alpha$: off $2\N$ the
density is nonzero (catastrophe), its derivative across $\alpha=2$ is
positive (tilt), and it vanishes on $2\N$ (oscillation). We record (iii) as
a theorem.

\begin{theorem}[Even moments carry no forced segment term]
\label{thm:dichotomy}
Let $m\ge1$ be an integer, $t\in\R$, and $h\in(0,\pi/8)$. Then the vertical
segment in \eqref{eq:split} contributes zero to
\[
  M_m(x;t):=\int_0^\infty u^{2m}g_t(u)\cos(xu)\,du,
\]
so that
\begin{equation}\label{eq:Mray}
  M_m(x;t)=\Re\int_0^\infty (s+ih)^{2m}\,g_t(s+ih)\,e^{ix(s+ih)}\,ds,
  \qquad
  |M_m(x;t)|\le C_m(t)\,e^{-hx},
\end{equation}
with $C_m(t)=\int_0^\infty|s+ih|^{2m}\,|g_t(s+ih)|\,ds<\infty$. Equivalently,
by the Erd\'elyi--Watson endpoint lemma of~\cite{Fsymbol}, the smooth even
weight $u^{2m}g_t(u)$ has an all-even-integer Taylor expansion at $u=0$,
whence its cosine transform is $O(x^{-K})$ for every $K$: $M_m$ has no
power tail. This is in exact contrast with the logarithmic weight, for
which the segment is the strictly positive Laplace transform
$N(x,t)=\tfrac\pi2\int_0^h y^2e^{-ty^2}\Phi(iy)e^{-xy}\,dy+O(e^{-hx})$
of~\cite{Fsymbol}.\footnote{The companion's endpoint asymptotics lemma
(Erd\'elyi--Watson form) states that a weight with an all-even-integer
Taylor scale at $u=0$ has cosine transform $O(x^{-K})$ for every $K$, every
even exponent carrying a vanishing cosine factor $\cos\tfrac{\pi(2j+1)}2=0$.}
\end{theorem}

\begin{proof}
By \eqref{eq:segment} with $w=u^{2m}$ and \cref{lem:density},
$\Im\,(iy)^{2m}=0$, so the segment real part is
$-\int_0^h0\cdot g_t(iy)e^{-xy}dy=0$. Hence $M_m(x;t)$ equals the real part
of the horizontal ray, which is \eqref{eq:Mray}; the bound is
$|e^{ix(s+ih)}|=e^{-hx}$ pulled out of the absolutely convergent integral
(finite because $g_t(s+ih)$ super-decays on the strip, with
$|\Phi(s+ih)|\le\sum_n(2\pi^2n^4e^{9s}+3\pi n^2e^{5s})
e^{-\pi n^2e^{4s}\cos4h}$ and $\cos4h>0$). The endpoint statement is the
companion's Erd\'elyi--Watson endpoint lemma: an all-even-integer Taylor
scale carries the vanishing cosine factor $\cos\tfrac{\pi(2j+1)}2=0$ at
every exponent.
\end{proof}

\Cref{thm:dichotomy} is the precise sense in which the even pseudo-process
directions are \emph{invisible to the positivity engine}. The contour shift
that makes the logarithmic response positive passes straight through the
even powers, leaving them to the oscillatory ray. We now show the ray
oscillates.

%==============================================================
\section{The moment-oscillation theorem}\label{sec:oscillation}
%==============================================================

The horizontal ray of \cref{thm:dichotomy} is an oscillatory integral whose
sign, at fixed $x_*\neq0$ and growing order $m$, is governed by a real
saddle that marches to infinity. We compute it on the real axis, where the
arithmetic of $\Phi$ is transparent.

\subsection{The arithmetic saddle}
For large $u$ the $n=1$, leading term of \eqref{eq:xi-rep} dominates:
\begin{equation}\label{eq:Phitail}
  \Phi(u)=2\pi^2e^{9u-\pi e^{4u}}\bigl(1+O(e^{-4u})+O(e^{-3\pi e^{4u}})\bigr),
  \qquad
  \log\Phi(u)=9u-\pi e^{4u}+\log2\pi^2+o(1),
\end{equation}
the $-3\pi e^{5u}$ correction and the $n\ge2$ terms being smaller by the
indicated super-exponential gaps. Write the amplitude exponent of
$M_m=\int_0^\infty e^{\rho_m(u)}e^{ix_*u}\,du$, with
$\rho_m(u):=2m\log u+tu^2+\log\Phi(u)$.

\begin{lemma}[Saddle location and curvature]\label{lem:saddle}
For each large $m$ there is a unique maximizer $u_m>0$ of $\rho_m$, with
\begin{equation}\label{eq:saddle}
  4\pi e^{4u_m}=\frac{2m}{u_m}+2tu_m+9,
  \qquad
  u_m=\tfrac14\log m-\tfrac14\log\log m+O(1)\to\infty,
\end{equation}
and $\rho_m''(u_m)=-16\pi e^{4u_m}+2t-\tfrac{2m}{u_m^2}
=-\dfrac{8m}{u_m}\bigl(1+o(1)\bigr)\to-\infty$, so the saddle is sharp with
width $\sigma_m:=|\rho_m''(u_m)|^{-1/2}=\sqrt{u_m/8m}\,(1+o(1))\to0$.
Moreover $u_m=u(m)$ is the restriction to integers of a function smooth in
real $m$, with $u'(m)=\tfrac1{4m}(1+o(1))$, so the increments
$u_{m+1}-u_m=\tfrac1{4m}(1+o(1))\to0$.
\end{lemma}

\begin{proof}
$\rho_m'(u)=\tfrac{2m}{u}+2tu+9-4\pi e^{4u}$ is strictly decreasing for
$u>0$ once $m$ is large (the $-16\pi e^{4u}$ in $\rho_m''$ dominates), from
$+\infty$ at $0^+$ to $-\infty$, giving a unique root $u_m$; that is the
first identity in \eqref{eq:saddle}. To leading order
$4\pi e^{4u_m}\sim 2m/u_m$, so $4u_m=\log m-\log(2\pi u_m)+o(1)$ and
$u_m=\tfrac14\log m-\tfrac14\log\log m+O(1)$. Substituting into
$\rho_m''$ and using $16\pi e^{4u_m}\sim 8m/u_m$ gives the curvature. For the
increment, the implicit function theorem applied to \eqref{eq:saddle} ---
whose $u$-derivative $16\pi e^{4u}+2m/u^2-2t$ is positive --- makes $u(m)$
smooth in real $m$; differentiating \eqref{eq:saddle} gives
$u'(m)=\bigl(2/u\bigr)\big/\bigl(16\pi e^{4u}+2m/u^2-2t\bigr)
=\tfrac1{4m}(1+o(1))$ by $16\pi e^{4u_m}\sim8m/u_m$, and the mean value
theorem yields $u_{m+1}-u_m=\tfrac1{4m}(1+o(1))\to0$. (Thus the $O(1)$ term
in \eqref{eq:saddle} has vanishing increments: it is a slowly varying
correction, not a bounded jitter.)
\end{proof}

\begin{lemma}[First-passage equidistribution]\label{lem:firstpassage}
Let $(\phi_m)_{m\ge1}\subset\R$ satisfy $\phi_m\to+\infty$ and
$\phi_{m+1}-\phi_m\to0$. Then $\{\phi_m\bmod2\pi\}$ meets every nonempty open
arc infinitely often; in particular $\cos\phi_m>\tfrac12$ for infinitely many
$m$, and $\cos\phi_m<-\tfrac12$ for infinitely many $m$. Genuine divergence
is essential: mere unboundedness of $\phi_m$ does not suffice, since a
trajectory with vanishing steps can turn back below an odd multiple of $\pi$
forever.
\end{lemma}

\begin{proof}
Fix a nonempty open arc; inside it take a closed sub-arc $J'$ of length
$\ell\in(0,2\pi)$, and choose $N$ with $|\phi_{m+1}-\phi_m|<\ell$ for all
$m\ge N$. Lift $J'$ to the disjoint real intervals $I_n=[a_n,a_n+\ell]$,
$a_{n+1}=a_n+2\pi$. For an arbitrary threshold $T$, divergence gives some
$m_0\ge N$ with $\phi_{m_0}>T+2\pi$; let $n$ be the largest index with
$a_n\le\phi_{m_0}$, so $a_n>T$. Put $k=\min\{j\ge m_0:\phi_j>a_n+\ell\}$,
finite since $\phi_j\to+\infty$. By minimality $\phi_{k-1}\le a_n+\ell$, and
$\phi_{k-1}\ge a_n$ --- else the single step $\phi_k-\phi_{k-1}>\ell$ would
contradict $|\phi_k-\phi_{k-1}|<\ell$ ($k-1\ge N$). Hence $\phi_{k-1}\in I_n$,
i.e.\ $\phi_{k-1}\bmod2\pi\in J'$ with $\phi_{k-1}\ge a_n>T$. As $T$ is
arbitrary, infinitely many such indices exist. The cosine statements follow
with the arcs $(-\tfrac\pi6,\tfrac\pi6)$ and $(\pi-\tfrac\pi6,\pi+\tfrac\pi6)$,
where $|\cos|>\tfrac{\sqrt3}2>\tfrac12$; if $\phi_m\to-\infty$ apply the
result to $-\phi_m$.
\end{proof}

\subsection{Oscillation}
\begin{theorem}[Moment oscillation]\label{thm:oscillation}
Fix $t\in\R$ and a real $x_*\neq0$. As $m\to\infty$,
\begin{equation}\label{eq:osc}
  M_m=\int_0^\infty u^{2m}e^{tu^2}\Phi(u)\cos(x_*u)\,du
  =A_m\cos\bigl(x_*u_m+\theta_m\bigr)\,\bigl(1+o(1)\bigr),
\end{equation}
with $A_m=\sqrt{2\pi}\,\sigma_m\,e^{\rho_m(u_m)}>0$ and a phase
$\theta_m=O(\sigma_m)\to0$. Consequently $M_m$ changes sign infinitely
often: both $\{m:M_m>0\}$ and $\{m:M_m<0\}$ are infinite, and the velocity
sign $\sgn(T_m'(0))=-\sgn(M_m/M_1)$ of \cref{thm:velocity} is not eventually
constant.
\end{theorem}

\begin{proof}
By \cref{lem:saddle} the mass of $e^{\rho_m(u)}$ concentrates in a window of
width $\sigma_m\to0$ about $u_m$, where the smooth phase $x_*u$ varies by
$O(x_*\sigma_m)=o(1)$. Laplace's method with a bounded analytic phase factor
gives
\[
  \int_0^\infty e^{\rho_m(u)}e^{ix_*u}\,du
  =\sqrt{2\pi}\,\sigma_m\,e^{\rho_m(u_m)}\,e^{i x_*u_m}\bigl(1+O(\sigma_m)\bigr),
\]
the $O(\sigma_m)$ collecting the cubic saddle correction
$\rho_m'''\sigma_m^3$ and the phase curvature $x_*^2\sigma_m^2$; the
contribution outside $|u-u_m|\le\sigma_m\log m$ is super-exponentially
smaller, and the $n\ge2$ and $e^{5u}$ corrections to \eqref{eq:Phitail} are
absorbed in $1+o(1)$. Taking real parts yields \eqref{eq:osc} with
$\theta_m=\arg(1+O(\sigma_m))=O(\sigma_m)\to0$. Finally set
$\phi_m:=x_*u_m+\theta_m$; we may take $x_*>0$, $\cos$ being even. By
\cref{lem:saddle} $\phi_m\to+\infty$ (genuine divergence, as
$u_m\to+\infty$), and
$\phi_{m+1}-\phi_m=x_*(u_{m+1}-u_m)+(\theta_{m+1}-\theta_m)
=\tfrac{x_*}{4m}(1+o(1))\to0$ (the second increment vanishing because
$\theta_m$ converges). \Cref{lem:firstpassage} then gives $\cos\phi_m>\tfrac12$
for infinitely many $m$ and $\cos\phi_m<-\tfrac12$ for infinitely many $m$;
since $A_m>0$ and the relative error in \eqref{eq:osc} is $o(1)$, it follows
that $M_m>0$ infinitely often and $M_m<0$ infinitely often.
\end{proof}

\begin{corollary}[Failure of the one-sided even-cone bridge]
\label{cor:nogo}
At every nondegenerate positive-$\Lambda$ collision $(x_*,\Lambda)$ with
$x_*\neq0$ (which holds always, since $H_\Lambda(0)=\int e^{\Lambda u^2}
\Phi>0$), the response-moment sequence $\{M_m\}_{m\ge1}$ satisfies:
\begin{enumerate}
\item[$(\mathrm i)$] it is \emph{not eventually one-signed}, so no cone
condition $\sgn(M_m)\equiv\text{const}$ holds and the one-sided variational
bridge fails;
\item[$(\mathrm{ii})$] it is \emph{not a Stieltjes moment sequence}: there is
no positive measure $\nu$ on $[0,\infty)$ with $M_m=\int w^m\,d\nu(w)$, since
such moments are nonnegative; so the response functional is not a positive
functional on the even generator cone, refuting the Hankel-positivity form
of the bridge.
\end{enumerate}
\end{corollary}

\begin{proof}
$(\mathrm i)$ is the last clause of \cref{thm:oscillation}. $(\mathrm{ii})$:
a Stieltjes moment sequence is nonnegative, contradicting $(\mathrm i)$.
\end{proof}

\Cref{cor:nogo} closes the one-sided even-cone bridge in its two clean forms
--- an eventual sign of $T_m'(0)$, equivalently of $\{M_m\}$, and a
positive-moment (Stieltjes) structure on $\{M_m\}$. It does not claim to
exclude every conceivable first-order mechanism, only the natural cone a
reader is led to try; through that cone the even pseudo-process hierarchy
does not reach RH.

%==============================================================
\section{The closure, and unification with the logarithmic route}
\label{sec:closure}
%==============================================================

The two first-order routes now read as one picture. Recall the response
measure $\mu_*$ of \eqref{eq:mustar} and the companion's clock function
\begin{equation}\label{eq:R}
  R(\beta):=\int_0^\infty u^{2\beta}\,d\mu_*
  =\frac{1}{D(x_*,\Lambda)}\int_0^\infty u^{2\beta}e^{\Lambda u^2}
   \Phi(u)\cos(x_*u)\,du,
\end{equation}
analytic in the half-plane $\Re\beta>-\tfrac12$, where the integral
converges: the super-exponential decay of $\Phi$ controls $u\to\infty$ for
every $\beta$, while at the endpoint $u\to0$ the factor $u^{2\beta}$ is
integrable against $\Phi(0)\neq0$ exactly when $\Re\beta>-\tfrac12$. As the
Mellin transform of the smooth even datum $e^{\Lambda u^2}\Phi(u)\cos(x_*u)$,
$R$ continues meromorphically to $\C$, with simple poles only at the
negative half-integers $\beta=-\tfrac12,-\tfrac32,\dots$ (the endpoint Taylor
scale), the residue at $-\tfrac12$ being $\Phi(0)/\bigl(2D(x_*,\Lambda)\bigr)$
--- the same endpoint constant that governs the catastrophe and the tilt.
Every value used below --- the germ at $\beta=1$, the integer samples
$m\ge1$, and the real zeros $\beta_k>0$ of \cref{sec:arithmetic} --- lies in
the half-plane of analyticity.
Its real samples are the data of both routes:
\begin{equation}\label{eq:samples}
  R(1)=1,\qquad
  R'(1)=2B_*=2\,\frac{N(x_*,\Lambda)}{D(x_*,\Lambda)},\qquad
  R(m)=\frac{M_m}{M_1}\ \ (m\in\N).
\end{equation}
Here $R(1)=1$ is the normalization $\int u^2d\mu_*=1$; $R'(1)=2\int u^2\log
u\,d\mu_*=2B_*$ by differentiating under the integral; and $R(m)=M_m/M_1$ is
\eqref{eq:velocity}.

\begin{theorem}[The two routes, unified]\label{thm:unified}
At a nondegenerate positive-$\Lambda$ collision the clock function $R$ of
\eqref{eq:R} satisfies:
\begin{enumerate}
\item[$(\mathrm i)$] \emph{Forced value and slope, from the
endpoint-singular directions.} $R(1)=1$ is pinned by time-rescaling, and
$R'(1)=2B_*>0$ is forced positive by the tilt positivity of~\cite{Fsymbol};
the latter refutes self-dual extremality, $T_*'(2)=-\Lambda B_*\neq0$.
\item[$(\mathrm{ii})$] \emph{Oscillating integer samples, from the
endpoint-smooth directions.} $R(m)=M_m/M_1$ changes sign infinitely often
(\cref{thm:oscillation}); no cone exists.
\end{enumerate}
The discriminant between $(\mathrm i)$ and $(\mathrm{ii})$ is the regularity
of the symbol at the endpoint $u=0$: singular weights ($\log$, fractional
powers) carry the imaginary-axis density $\kappa$ and a forced sign, smooth
even weights sit at the zeros of $\kappa$ and oscillate. Both first-order
routes to $\Lambda\le0$ therefore close, for complementary reasons.
\end{theorem}

\begin{proof}
$(\mathrm i)$ is \eqref{eq:samples} together with the companion's tilt
theorem ($N>0$ on the admissible range) and collision-velocity formula.
$(\mathrm{ii})$ is \cref{thm:oscillation}. The discriminant statement is
\cref{thm:dichotomy} and the three-phenomena reading of
\cref{sec:dichotomy}.
\end{proof}

\begin{remark}[Why the slope is forced but the samples are free]
\label{rem:freevalue}
The companion's clock proposition already isolated the structural fact:
inside the subordination strip the clock sector fixes $R(1)=1$ and
\emph{leaves $R'(1)$ free}. \Cref{thm:unified} fills in both ends with the
arithmetic of $\Phi$. The slope $R'(1)=2B_*$ is forced by the one singular
tangent in the strip, where the contour segment is active; the values
$R(m)$, $m\ge2$, are free of any segment and realize their freedom by
oscillating. The companion was pointing at the right place --- the content
sits in the derivative, not the value --- and the present theorem shows the
even hierarchy adds no new pinned quantity: it samples the free part of $R$,
and the free part has no fixed sign.
\end{remark}

The function $R$ is thus an object with a forced germ at $\beta=1$
($R(1)=1$, $R'(1)>0$) and oscillating integer values --- perfectly possible
for a function analytic in $\Re\beta>-\tfrac12$, and exactly what occurs.
Any surviving route must
constrain $R$ \emph{globally}, not through the sign of any one sample; this
is taken up in \cref{sec:frontier}.

%==============================================================
\section{The probabilistic reading: what the level-$p$ calculus contributes}
\label{sec:probabilistic}
%==============================================================

The even members of \eqref{eq:family} are not an arbitrary test family: for
$\alpha=2m$ the multiplier $e^{tu^{2m}}$ is the symbol of the order-$2m$
heat pseudo-process of Krylov and Hochberg~\cite{Krylov,Hochberg}, run in
the \emph{backward} (growing) direction that is the de Bruijn--Newman
direction. The companion paper~\cite{Fpseudo} builds the representation
calculus for these pseudo-processes. We read the present result through it,
and find that it supplies the language of the response measure and,
decisively, the reason the bridge cannot succeed.

\subsection{The pseudo-process behind $\alpha=2m$}
The order-$p$ pseudo-process $X^{(p)}$ has characteristic functional
$\E[e^{ikX^{(p)}_t}]=e^{c_p(ik)^pt}$ and signed transition kernel $K_p$
(Gaussian at $p=2$, Airy at $p=3$, Hochberg's biharmonic at
$p=4$)~\cite{Fpseudo}. For $p=2m$ even the Fourier symbol is
$e^{c_{2m}(-1)^mk^{2m}t}$; matching the de Bruijn--Newman weight $e^{tu^{2m}}$
on the cosine side ($k\leftrightarrow u$, sign $c_{2m}(-1)^m=+1$) identifies
$H_{t,2m}$ as $H_0$ evolved by the order-$2m$ pseudo-process semigroup in
its backward, non-contractive direction --- the direction outside the
probabilistic strip, where~\cite{Fsymbol} already located the catastrophe
and~\cite{Fpseudo} the loss of every positive structure.

\subsection{The representer dictionary}
The response formula $\frac{d}{d\epsilon}T_*|_\sigma=-\Lambda\int\sigma\,
d\mu_*$ of \eqref{eq:mustar} is a Clark--Ocone pairing in the precise sense
of~\cite{Fpseudo}: the derivative of the observable $T_*$ along a direction
$\sigma$ is the pairing of $\sigma$ against one fixed representer $\mu_*$,
exactly as a centered functional of a noise pairs against its
representer, and the order-$p$ calculus is what makes ``derivative defined
by adjointness against a divergence, with no positivity'' a theorem rather
than an analogy. In that dictionary the extremality hypothesis along the
stable tangent is a first-chaos-type orthogonality
$B_*=\int u^2\log u\,d\mu_*=0$, and the even-hierarchy cone is the family of
samples $R(m)=\int u^{2m}d\mu_*$. The companion~\cite{Fpseudo} gives these
pairings a rigorous home; it does not, by itself, decide their signs.

\subsection{Why the calculus forecloses the cone rather than enabling it}
The decisive contribution of~\cite{Fpseudo} is negative and sharp. The cone
the bridge needs is a \emph{positivity} statement about the even generator
directions; the orders $p=2m\ge4$ that would have to supply it are exactly
the orders from which positivity has been stripped. Three facts
of~\cite{Fpseudo} bear directly:
\begin{enumerate}
\item[$(\mathrm i)$] For $p\ge3$ there is no positive path-space law: the
pseudo-pairing is signed, $\E[X_T^2]=0$, and $X^{(p)}$ is not a
semimartingale. There is no positive measure whose moments are the responses,
which is \cref{cor:nogo}$(\mathrm{ii})$ seen from the process side.
\item[$(\mathrm{ii})$] Orthogonal chaos exists if and only if $p=2$
(\cite{Fpseudo}, orthogonality theorem): the Appell--Gram matrix is diagonal
only in the Gaussian case. There is no orthogonal first chaos for the even
direction to be orthogonal \emph{to}; only a triangular, biorthogonal
structure survives, which carries no sign.
\item[$(\mathrm{iii})$] Correspondingly --- and this is the companion's
even-exponent rigidity~\cite{Fsymbol} --- no $e^{\epsilon u^{2m}}$,
$m\ge2$, is a P\'olya universal factor: the Laguerre--P\'olya preservation
that real-rootedness needs exists only at $p=2$. A direction that does not
preserve real-rootedness cannot point uniformly into the real-rooted side,
so its collision response has no reason to be one-signed --- and is not.
\end{enumerate}

\subsection{The separation of the two moment systems}
The cleanest statement of why the calculus cannot transfer is that the
pseudo-process carries \emph{two different moment systems}, and RH needs the
one the calculus does not control. The process's own moments are
sign-regular and explicit,
\begin{equation}\label{eq:procmoments}
  \E\bigl[(X^{(p)}_T)^{pn}\bigr]=\frac{(pn)!}{n!}\,c_p^{\,n}\,T^n,
  \qquad \E\bigl[(X^{(p)}_T)^{m}\bigr]=0\ \ (p\nmid m),
\end{equation}
a single sign $\sgn(c_p^n)$, governed by the grading $p\mid m$
of~\cite{Fpseudo}. The \emph{arithmetic} response moments, by contrast, are
\begin{equation}\label{eq:arithmoments}
  M_m=\int_0^\infty u^{2m}\,e^{\Lambda u^2}\Phi(u)\cos(x_*u)\,du,
\end{equation}
moments of the \emph{signed} measure $\mu_*$, whose sign is twisted by the
zero's phase $\cos(x_*u)$. The level-$p$ Clark--Ocone calculus governs
\eqref{eq:procmoments}: the grading, the representation, the Stein identity,
the carr\'e du champ. It is silent on \eqref{eq:arithmoments}, because the
$\cos(x_*u)$ factor is not part of the pseudo-process data but of the
\emph{arithmetic} of $\Phi$ at the collision. The two systems agree only in
the trivial sector; the calculus cannot lend its sign-regularity to the
oscillating one. That silence is the precise sense in which the
pseudo-process framework, however complete as probability, does not open the
Riemann hypothesis.

\begin{remark}[The honest summary of {\cite{Fpseudo}}'s role]
\label{rem:honest}
The level-$p$ Clark--Ocone paper helps in two ways and not in a third. It
gives the rigorous backbone for the response-measure / representer language
used throughout~\cite{Fsymbol} and here (the adjoint definition of the
operator derivative, the uniqueness of the representer up to chronological
equivalence). And it diagnoses the obstruction: the cone is a positivity
condition, and the orders that would carry it are positivity-free. What it
does not do --- what nothing in the pseudo-process category does --- is
constrain the signed arithmetic moments $M_m$; those are decided by the
contour mechanism of \cref{sec:dichotomy}, which lives in the analysis of
$\Phi$, not in the probability of $X^{(p)}$.
\end{remark}

%==============================================================
\section{Numerical confirmation}\label{sec:numerics}
%==============================================================

We confirm the oscillation directly on the Riemann kernel. The integrals
$M_m(x_0;t)=\int_0^\infty u^{2m}e^{tu^2}\Phi(u)\cos(x_0u)\,du$ are evaluated
by subdivided tanh--sinh quadrature on $u\in[0,2.1]$ at $50$-digit working
precision, $\Phi$ summed until the terms underflow. Two checks pin the
conventions: $\Phi(0)=0.44669690$, matching a direct $30$-digit summation of
$\sum_{n\ge1}(2\pi^2n^4-3\pi n^2)e^{-\pi n^2}=0.4466969004671234$, and
$H_0(2\gamma_1)=\int\Phi\cos(2\gamma_1u)\,du=-4.6\times10^{-14}$, vanishing
to the working precision, where $\gamma_1=14.134725$ is the first zeta
ordinate. Since the collision site of a hypothetical positive-$\Lambda$
scenario is not known (its existence is RH), we exhibit the mechanism at the
real surrogate sites $x_0=2\gamma_j$, where the oscillation is generic in
$x_0\neq0$; \cref{thm:oscillation} applies verbatim at any genuine
collision.

\begin{center}
\begin{tabular}{rrcrc}
\toprule
$m$ & $M_m(x_0;0)$ & $\sgn M_m$ & $-M_m/M_1$ & $\sgn T_m'$\\
\midrule
1  & $-5.009\times10^{-5}$ & $-$ & $-1.000$            & $-$\\
2  & $\phantom{-}3.709\times10^{-6}$ & $+$ & $\phantom{-}0.0740$ & $+$\\
3  & $\phantom{-}1.752\times10^{-7}$ & $+$ & $\phantom{-}0.00350$ & $+$\\
4  & $-6.033\times10^{-10}$ & $-$ & $-1.20\times10^{-5}$ & $-$\\
5  & $-1.119\times10^{-9}$ & $-$ & $-2.24\times10^{-5}$ & $-$\\
$\vdots$ & $\vdots$ & $-$ & $\vdots$ & $-$\\
9  & $-1.144\times10^{-13}$ & $-$ & $-2.28\times10^{-9}$ & $-$\\
10 & $\phantom{-}5.151\times10^{-15}$ & $+$ & $\phantom{-}1.03\times10^{-10}$ & $+$\\
$\vdots$ & $\vdots$ & $+$ & $\vdots$ & $+$\\
20 & $\phantom{-}1.376\times10^{-21}$ & $+$ & $\phantom{-}2.75\times10^{-17}$ & $+$\\
21 & $\phantom{-}4.976\times10^{-25}$ & $+$ & $\phantom{-}9.93\times10^{-21}$ & $+$\\
22 & $-8.387\times10^{-23}$ & $-$ & $-1.67\times10^{-18}$ & $-$\\
$\vdots$ & $\vdots$ & $-$ & $\vdots$ & $-$\\
30 & $-1.840\times10^{-26}$ & $-$ & $-3.67\times10^{-22}$ & $-$\\
\bottomrule
\end{tabular}
\end{center}

\noindent At $x_0=2\gamma_1=28.2694$, $t=0$, the sign pattern of $M_m$ over
$m=1,\dots,30$ is
\[
  -\,+\,+\,-\,-\,-\,-\,-\,-\,+\,+\,+\,+\,+\,+\,+\,+\,+\,+\,+\,+\,
  -\,-\,-\,-\,-\,-\,-\,-\,-,
\]
with \textbf{four sign reversals}; the near-vanishing $M_{21}=4.98\times
10^{-25}$, three orders below its neighbours, is the amplitude dip of a
cosine crossing zero, an incipient fifth reversal. Three readings confirm
the theorem. First, the magnitudes fall geometrically, $|M_m|$ shrinking
by roughly an order per step, as the saddle amplitude
$A_m=\sqrt{2\pi}\sigma_me^{\rho_m(u_m)}$ requires (the weight $u^{2m}$ at the
sub-unit saddle $u_m<1$ contracts the integral). Second, the oscillation is
\emph{faster at larger $x_0$}: at $x_0=2\gamma_3=50.0217$ the same sweep
reverses sign \textbf{six} times, matching the saddle spacing
$\Delta(\log m)=4\pi/x_0$ between reversals (a factor $e^{4\pi/x_0}$ in $m$),
which predicts the ratio $50/28\approx1.8$ in reversal counts against the
observed $6/4=1.5$, the discrepancy being the slow onset of the asymptotic
at small $m$. Third, the pattern is \emph{stable under the flow}: repeating
at $t=0.2$ leaves the sign sequence and reversal count unchanged
($-{+}{+}{-}\cdots$, four reversals), as \cref{thm:oscillation} requires,
the leading constant being $t$-independent. The velocity sign
$\sgn(T_m'(0))=-\sgn(M_m/M_1)$ inherits the reversals exactly: the even
pseudo-process directions push the collision time alternately forward and
backward, and no cone exists. Simulation code accompanies the source.

%==============================================================
\section{The arithmetic of the clock function}\label{sec:arithmetic}
%==============================================================

The frontier object is the clock function $R(\beta)=\Mcal(\beta)/\Mcal(1)$
of \eqref{eq:Mcal}. It is not a generic analytic function: at the critical
value $\Lambda=0$ it is the even-derivative family of the completed zeta
function at the collision height, so its samples are arithmetic and its real
zeros are a log-periodic spectrum whose period is the collision location.

\subsection{The exact zeta identity}
Recall the classical $H_0(x)=\tfrac18\,\Xi(x/2)$, with $\Xi(t)=\xi(\tfrac12+
it)$ the Riemann $\Xi$-function, real and even on $\R$.

\begin{theorem}[The clock samples are even derivatives of $\Xi$]
\label{thm:xi}
For $\Lambda=0$, every integer $m\ge0$, and every real $x_*$,
\begin{equation}\label{eq:xi-id}
  M_m=\int_0^\infty u^{2m}\Phi(u)\cos(x_*u)\,du
     =(-1)^m H_0^{(2m)}(x_*)
     =\frac{(-1)^m}{8\cdot4^m}\,\Xi^{(2m)}\!\Bigl(\frac{x_*}{2}\Bigr).
\end{equation}
In particular $\Mcal(0)=\tfrac18\Xi(x_*/2)$, which vanishes precisely when
$\tfrac12+i\tfrac{x_*}2$ is a nontrivial zero of $\zeta$. For $\Lambda>0$ the
same identities hold with $H_0$ replaced by the heat-flowed $H_\Lambda$; at
a genuine collision $\Mcal(0)=H_\Lambda(x_*)=0$ automatically, $x_*$ being a
zero of $H_\Lambda$.
\end{theorem}

\begin{proof}
Differentiating $H_0(x)=\int_0^\infty\Phi(u)\cos(xu)\,du$ under the integral
and using $\tfrac{d^{2m}}{dx^{2m}}\cos(xu)=(-1)^mu^{2m}\cos(xu)$ gives
$H_0^{(2m)}(x_*)=(-1)^mM_m$; the chain rule on $H_0(x)=\tfrac18\Xi(x/2)$ gives
$H_0^{(2m)}(x)=\tfrac1{8\cdot4^m}\Xi^{(2m)}(x/2)$. The $\zeta$-zero statement
is the classical representation. The $\Lambda>0$ version replaces $\Phi$ by
$e^{\Lambda u^2}\Phi$ and $H_0$ by $H_\Lambda$.
\end{proof}

Identity \eqref{eq:xi-id} was confirmed to $35$ digits by evaluating both
sides independently --- the moment integral over $\Phi$, and the Taylor
derivatives of $\Xi$ obtained from $\zeta$ --- at the zero height
$x_*=2\gamma_1$ (where $\Mcal(0)$ and $\Xi(\gamma_1)$ both vanish to
$10^{-14}$) and at a generic height. This identity, which the Mellin
structure of $\Phi$ makes natural, is exact: the clock function is the
moment/derivative family of completed zeta at the collision height, and the
continuous-$\beta$ clock is its fractional even-derivative interpolation. The Dirichlet-polynomial structure --- each
$n$-term of $\Phi$ contributing the character $n^{-1/2-ix_*/2}$ after the
Mellin exchange $v=e^{4u}$ --- is exactly the Riemann--Siegel expansion of
the $\zeta(\tfrac12+i\tfrac{x_*}2)$ inside $\Xi$, read through the
$\beta$-derivatives, with the classical cutoff $n\lesssim\sqrt{\beta/2\pi}$
furnished by the per-term saddle; it is zeta's own approximate functional
equation, not a new one.

\begin{corollary}[The oscillation is a P\'olya--Jensen sign failure]
\label{cor:jensen}
\Cref{thm:oscillation} is the statement that the even derivatives
$\Xi^{(2m)}(x_*/2)$ do not settle to the alternating sign $(-1)^m$. The cone
condition $\sgn(M_m)\equiv\text{const}$ is the sign-regularity of the
P\'olya--Jensen sequence $\bigl\{(-1)^m\Xi^{(2m)}(x_*/2)\bigr\}_m$ shifted to
height $x_*/2$. At the central height $x_*=0$ this sequence is of one sign,
$M_m(0)=\int u^{2m}\Phi>0$, the Laguerre--P\'olya positivity whose Jensen
polynomials Griffin--Ono--Rolen--Zagier prove eventually
hyperbolic~\cite{GORZ}; \cref{thm:oscillation} shows the sign-regularity
fails at every other height. The cone holds exactly at $x_*=0$, which is
never a collision ($H_\Lambda(0)>0$). This is why every heuristic in the
program degenerated ``unless $x_*=0$'': the central height is the one place
the P\'olya--Jensen sequence is monotone in sign.
\end{corollary}

\subsection{The frequency law and approximate discrete scale invariance}
The real zeros of $R$ inherit the saddle of \cref{thm:oscillation}: a zero
needs $x_*u_\beta=\tfrac\pi2+k\pi$ with $u_\beta=\tfrac14\log\beta-\tfrac14
\log\log\beta+O(1)$, whence
\begin{equation}\label{eq:freq}
  \log\beta_k=\frac{4\pi}{x_*}\,k+\log k+O(1).
\end{equation}
The real zeros form an asymptotic geometric progression of ratio
$e^{4\pi/x_*}$: a log-periodic (discrete-scale-invariant) spectrum whose
log-period $4\pi/x_*$ is the collision location read as a frequency. The
progression is only asymptotic --- the drift $\beta_{k+1}/\beta_k=
e^{4\pi/x_*}(1+\tfrac1k+o(\tfrac1k))$ decays like $1/k$. Numerically, the
sign-change gaps of $M_m$ at $x_*=2\gamma_1$ are $0.92,0.79,0.65,0.63$ across
$m\le120$, matching the second-order law $4\pi/x_*+\tfrac1k=0.94,0.78,0.69,
0.64$ at $k=2,\dots,5$: the data confirm the $\log k$ correction of
\eqref{eq:freq}, not merely its limit $4\pi/x_*=0.444$ (the first period,
$k=1$, with zeros near $\beta=1.2,4.0$, is pre-asymptotic). At $x_*=2\gamma_3$
ten periods track the same law toward $0.251$. The
breaking of exact scale invariance is the arithmetic of \cref{thm:xi}:
discrete scale invariance is the leading shadow of the zeta derivatives, and
the dilation generator whose spectrum this is, is the Mellin dilation
itself, acting on the auxiliary $\beta$-plane, not a Hilbert--P\'olya
operator on the zeros of $\zeta$.

\subsection{Self-reference, and why it does not propagate}
By \cref{thm:xi} the order-$m$ clock sample, as the height $x_*$ varies,
traces $\Xi^{(2m)}(x_*/2)$. The zero heights are marked \emph{exactly} at
$m=0$, where $\Mcal(0)=\tfrac18\Xi$ vanishes at every ordinate; but the
marking does not survive differentiation. Sweeping $x_*$ across
$2\gamma_1=28.269$ and $2\gamma_2=42.044$: at $\beta=0$ the zeros of $M_0$
sit on the ordinates; at $\beta=3$ the only zero of $M_3$ in $[26,45]$ lies
$3.5$ and $10.3$ away from them; at $\beta=10$ the nearest zero of $M_{10}$
misses $2\gamma_1$ by $0.23$ and $2\gamma_2$ by $1.9$. The zero-height
signature is exact at the bottom of the ladder and effectively gone within a
few rungs: the high even moments that defeat the cone are decoupled from the
\emph{locations} of the zeros, even though they are built from $\Xi$. This is
the differentiation form of the decoupling that \cref{sec:probabilistic}
reads probabilistically; the response is arithmetic, but the arithmetic that
survives to high order is the local Taylor data of $\Xi$, not its zero set.

\subsection{The quartic flow and the decoherence of the zero field}
\label{subsec:quartic}
The self-reference asked how far \emph{up in $\beta$} the zero-coherence of
the clock survives; the companion's quartic flow $H_{t,4}$ poses the partner
question, how far \emph{up in $x$} the interlacing of the zeros survives
under averaging. It is the one even member inside the probabilistic strip:
the iterated-Brownian (Funaki) average $H_{t,4}=\E_{s\sim N(0,2t)}[H_s]$ of
the de Bruijn--Newman flow over a Gaussian of signed times, validated here
to $13$--$15$ digits. Every nonreal zero of $H_{t,4}$ is therefore a
cancellation in an average of mixands each of which (for $s\ge\Lambda$) is
real-rooted --- the average destroying a realness that each term possesses.
Fine-grid tracking exhibits it directly. Under the flow the pair of zeros at
$2\gamma_4=60.85$ and $2\gamma_5=65.87$ closes (to $62.19,\,63.81$ at
$t=6$), collides near $x\approx63$, and lifts off: the value $H_{t,4}(63)$
between them passes from a near-zero touching ($8\times10^{-13}$ at $t=6$,
two real zeros) to a non-crossing dip that recedes from the axis
($5\times10^{-11}$ at $t=7$, $1.7\times10^{-10}$ at $t=8$, no real zero) ---
the sign-indefinite unfolding of the companion resolving to the nonreal
side. The pair stays nonreal through $t=10$, so numerically
\begin{equation}\label{eq:lam4}
  \Lambda(4)>10,
\end{equation}
a first quantitative lower bound; the proven window-realness of the
companion is thus a large-$t$ \emph{re}-coherence, absent at moderate times.

The pre-registered coherence criterion --- decoherence when the de
Bruijn--Newman zero displacement across the window, $|v_j|\sqrt{2t}$ with
$v_j=\partial_sx_j=H_0''/H_0'$, reaches the half-gap $h_j$, i.e.\ at
$t_c(x_j)=h_j^2/2v_j^2$ --- assigns the \emph{smallest} coherence time to
exactly this pair ($t_c=3.45$ at $2\gamma_4$, the minimum over the low
zeros) and predicts the onset to a factor of two ($3.45$ against
$\approx6.5$). To that extent the interlacing lens is right: it names the
first-colliding pair. But the velocities $|v_j|$ are $O(1)$ and fluctuating
(mean $0.66$ over $\gamma\le95$, no monotone trend), so there is no smooth
polynomial $x_c(t)$; the decoherence is configuration-specific, set by the
local minimum of $h_j^2/v_j^2$, and the nucleation height $x\approx63$ sits
near the upper (saddle) edge $x\sim t\log^3t$ of the companion's band at
$t\approx6.5$, not its lower edge. The lens predicts which pair, and roughly
when, but not a scaling law.

%==============================================================
\section{What survives: the global frontier}\label{sec:frontier}
%==============================================================

The first-order even bridge is closed, but the closure is sharp enough to
say where an opening, if one exists, must now live. It cannot be a pointwise
sign of a response; those oscillate. It must be a \emph{global} constraint
on the function
\begin{equation}\label{eq:Mcal}
  \Mcal(\beta):=\int_0^\infty u^{2\beta}\,e^{\Lambda u^2}\Phi(u)
  \cos(x_*u)\,du,
  \qquad R(\beta)=\Mcal(\beta)/\Mcal(1),
\end{equation}
whose germ at $\beta=1$ is forced ($R(1)=1$, $R'(1)=2B_*>0$) and whose
integer samples oscillate. Three problems sharpen this.

\begin{problem}[The Jensen-sequence form of the global bridge]
\Cref{sec:arithmetic} identifies $\Mcal$ exactly: its integer samples are
the shifted P\'olya--Jensen sequence $(-1)^m\Xi^{(2m)}(x_*/2)$
(\cref{thm:xi}), and its real zeros are the log-periodic spectrum
\eqref{eq:freq}. The global bridge is therefore not a question about generic
analytic functions but the sharp arithmetic question: can the
P\'olya--Jensen sequence of $\Xi$ at a height $x_*/2$ be sign-regular when
$x_*$ is the site of a real multiplicity-two zero of the flowed $\Xi$? By
\cref{cor:jensen} sign-regularity holds only at the central height $x_*=0$,
which is never a collision; the open content is whether the collision
constraint (a double zero of $H_\Lambda$ at $x_*$, with the
Laguerre--P\'olya structure of $H_\Lambda$ in force) can coexist with the
oscillation forced by \cref{thm:oscillation}. Because the relevant
$\beta$-aspect data is zeta's own Dirichlet polynomial (\cref{thm:xi}), the
tools of the Dirichlet-polynomial genre apply --- Tur\'an power sums,
zero-density estimates --- but to the same hard object as $\zeta$ itself: the
reframing arithmetizes the bridge without escaping zeta's difficulty, and is
self-referential, the response spectrum at $x_*$ being built from zeta data
at $x_*$. Unlike the first-order signs, it is not a pointwise positivity and
not obviously doomed.

\smallskip\noindent\emph{Partial resolution (an $\mathcal{LP}$ reformulation).}
Since $x_*$ is a double zero of $H_\Lambda\in\mathcal{LP}$, write
$H_\Lambda(x_*+s)=s^2h(s)$ with $h\in\mathcal{LP}$ and
$h(0)=\tfrac12H_\Lambda''(x_*)\neq0$ (removing a real double zero preserves
the Laguerre--P\'olya class). Then
\begin{equation}\label{eq:LPform}
  M_m=(-1)^m H_\Lambda^{(2m)}(x_*)=(-1)^m(2m)(2m-1)\,h^{(2m-2)}(0),
\end{equation}
so the cone $\sgn(M_m)\equiv\text{const}$ is exactly sign-regularity of the
even Taylor data of $h$. For a \emph{one-sided} zero set --- e.g.\ $h$ of
genus zero with all zeros negative, $h(s)=h(0)\prod_k(1+s/a_k)$, $a_k>0$ ---
every $h^{(j)}(0)>0$ and the cone holds; the collision's zeros, the remaining
zeros of $H_\Lambda$ shifted to straddle the origin, are two-sided, and
\cref{thm:oscillation} shows that two-sidedness breaks the regularity. This
recovers the oscillation inside $\mathcal{LP}$ and locates its cause in the
\emph{bulk} position of the collision; it is structurally consistent with
$\Lambda>0$, so the obstruction, if any, is in the global zero geometry, not
the local Jensen data.
\end{problem}

\begin{problem}[A second invariant beyond first order]
The velocity law $T_*'(2)=-\Lambda B_*$ fixes the corner slope of $R$; the
even samples fix nothing. Following the companion's open
problem~\cite{Fsymbol}, seek a \emph{second} collision invariant, forced in
every positive-$\Lambda$ scenario and incompatible with both the
nondegenerate velocity $T_*'(2)\neq0$ and the degenerate cusp. The natural
candidates are the fractional clock responses $\mathcal G_\beta H$ at the
collision (the companion's clock calculus), among which a forced relation
would be a constraint a positive-$\Lambda$ scenario must satisfy; the even
hierarchy supplies the integer-$\beta$ values of this family, now known to
oscillate, which constrains the interpolating $R$.

\smallskip\noindent\emph{Partial resolution (the forced germ).} The contour
engine of \cref{thm:dichotomy} pins not only $R'(1)$ but the entire germ of
$R$ at $\beta=1$. Each log-moment
$N_k=\int_0^\infty u^2(\log u)^k e^{\Lambda u^2}\Phi\cos(x_*u)\,du$ is
endpoint-dominated, its segment density $-\Im[(iy)^2\Log^k(iy)]$ giving
$N_k/N_1\sim k\,G^{(k-1)}(3)/G(3)$ with $G(s)=\Gamma(s)x_*^{-s}$; e.g.\
$N_2/N_1\sim2(\psi(3)-\log x_*)$, confirmed numerically to under one percent
by $x_*\ge120$ and exactly at $x_*=480$. Summing the tower
($R^{(k)}(1)/R'(1)=2^{k-1}k\,G^{(k-1)}(3)/G(3)$) resums to a closed form,
\begin{equation}\label{eq:germ}
  R(\beta)=1+B_*\,(\beta-1)\,\Gamma(2\beta+1)\,x_*^{-2(\beta-1)}
  \;+\;(\text{beyond all orders}),
\end{equation}
a $\Gamma$-factor times the scale $x_*^{-2(\beta-1)}$ --- the $\Gamma$-factor
coefficient anticipated by the $\beta$-aspect reading of \cref{thm:xi}. So the
second invariant, and indeed a whole tower, \emph{is} found, in closed form;
but it is determined by the two endpoint numbers $B_*$ and $x_*$ alone and is
consistent with every $\Lambda>0$. The RH-relevant content is not in the
forced germ but in the beyond-all-orders remainder --- the oscillatory zeros
$\beta_k$ of \eqref{eq:freq} --- exactly where the arithmetic of
\cref{sec:arithmetic} lives.
\end{problem}

\begin{problem}[The quartic flow: the one positive even direction]
\label{prob:quartic}
Among the even exponents only $\alpha=4$ lies in the probabilistic
subordination strip $[2,4]$: the quartic backward flow is the
iterated-Brownian (Funaki) evolution of $\Xi$, a Gaussian average of the de
Bruijn--Newman trajectory at a random signed time~\cite{Fsymbol,Funaki},
and it is the one even pseudo-process that retains a positive clock. The
companion proves $H_{t,4}$ real-rooted on every fixed window for large $t$;
whether $\Lambda(4)<\infty$ globally is open, and \cref{subsec:quartic} now
bounds it from below, $\Lambda(4)>10$, by exhibiting a bulk collision that
lifts a conjugate pair off the axis at $t\approx6.5$. This is the surviving
frontier where positivity is not lost: a global real-rootedness theorem for
the quartic flow would constrain the de Bruijn--Newman flow underneath it
through the Gaussian average, a mechanism the sign-oscillating odd-order
even powers cannot provide. The even hierarchy fails as a \emph{cone};
whether its single positive member fares better as an \emph{average} is
settled below, and it does not.

\smallskip\noindent\emph{Partial resolution (toward $\Lambda(4)=\infty$).}
Tracking the lifted pair of \cref{subsec:quartic} to larger $t$ sharpens
\eqref{eq:lam4}: its normalized axis distance
$\min_{[58,68]}|H_{t,4}|/\max_{[58,68]}|H_{t,4}|$ holds near $0.02$--$0.04$
through $t=300$ with no sign of return, so $\Lambda(4)>300$. But more than the
bound, the geometry tells against re-coherence ever winning. Two scales run
opposite ways. A pair re-coheres only when the Laplace window $x\sim t^{1/4}$
climbs to its height --- cost $t\sim x^4$, polynomial and slow. Decoherence,
by contrast, reaches \emph{up}: the linear criterion of \cref{subsec:quartic}
gives onset $t_c(x)\sim h(x)^2/2v^2$, and as the height grows the zero gaps
close, $h(x)\sim2\pi/\log x\to0$, while the velocities stay $O(1)$, so
$t_c(x)\to0$. Hence at every fixed $t>0$ there are heights --- all
sufficiently large $x$ --- whose pairs have already lifted off, and the
decohered region descends from $x=\infty$ as fast as the gaps close while the
real-rooted window rises only as $t^{1/4}$. The band between them never
empties, so $H_{t,4}$ carries nonreal zeros at every $t$, and conjecturally
\begin{equation}\label{eq:lam4inf}
  \Lambda(4)=\infty .
\end{equation}
This is consistent with the companion's fixed-window realness --- that window
is exactly the climbing $t^{1/4}$ edge --- and is, we expect, provable: it
needs only an effective Laplace lower edge and a gap-compression estimate
exhibiting a genuine lifted pair above it at each $t$. The quartic average
does not rescue real-rootedness; its positivity is precisely what fills the
dips and lifts the pairs.
\end{problem}

\subsection{The global object, and where the barrier sits}
\label{subsec:global}
The three partial resolutions agree: the obstruction to $\Lambda\le0$, if it
exists, is in no local or forced datum --- not the germ \eqref{eq:germ}, the
local Jensen sequence \eqref{eq:LPform}, or the finite-$t$ quartic --- but in
the global structure. That structure has an exact name. Suppose $\Lambda>0$
with a nondegenerate collision $(x_*,\Lambda)$, and define the \emph{collision
trajectory}
\begin{equation}\label{eq:traj}
  \varphi(s):=\frac{H_{\Lambda+s}(x_*)}{D},
  \qquad D=D(x_*,\Lambda)=-H_\Lambda''(x_*),
\end{equation}
the de Bruijn--Newman flow read at the fixed collision height as a function
of flow time. It is entire in $s$, with $\varphi(0)=0$ (the collision is a
zero) and $\varphi'(0)=\partial_tH_\Lambda(x_*)/D=1$. Since $H_\Lambda(x_*)=0$,
the clock calculus of~\cite{Fsymbol} makes the entire clock function the
Marchaud fractional derivative of $\varphi$ from the pre-threshold side:
\begin{equation}\label{eq:globalclock}
  R(\beta)=\frac{\beta(\beta-1)}{\Gamma(2-\beta)}\int_0^\infty
  \bigl[\varphi(-r)+r\bigr]\,r^{-1-\beta}\,dr,\qquad 1<\beta<2,
\end{equation}
the compensator $+r$ converging because $\varphi(-r)+r=O(r^2)$. The whole
program collapses onto this one object: $R$ is an order-$\beta$ fractional
derivative of $\varphi$; its forced germ \eqref{eq:germ} is $\varphi$ near
$s=0$; its log-periodic zeros \eqref{eq:freq} are the oscillation of
$\varphi(-r)$ in $r$. The Riemann hypothesis is the statement that no such
trajectory exists with $\Lambda>0$.

This is a faithful reduction, and it shows precisely why it is not an escape.
The trajectory is not free. At $s=-\Lambda$,
\begin{equation}\label{eq:barrier}
  \varphi(-\Lambda)=\frac{H_0(x_*)}{D}=\frac{\Xi(x_*/2)}{8D},
\end{equation}
the completed zeta function at the collision height, and for $s<-\Lambda$ the
backward continuation of the flow past $t=0$. The clock $R$ therefore encodes
the flow trajectory \emph{through} the value $\Xi(x_*/2)$: the global object
is zeta itself, read through the heat flow. The local shots showed the
obstruction is not in the germ; the global object shows that the germ's
complement is the arithmetic. We have localized where $\Lambda\le0$ must be
decided --- the Mellin/fractional structure of the single trajectory
$\varphi$, equivalently the oscillation of the backward flow at the collision
height --- without dissolving it, because that trajectory is the Riemann
$\Xi$-function in motion. This is the honest endpoint of the deformation
program: a complete map of the obstruction's location, and a barrier exactly
as hard as the hypothesis.

\begin{remark}[Even-data blindness: why every route meets the same barrier]
\label{rem:blind}
There is a single structural reason behind all of the above. Since
$\partial_t=-\partial_x^2$, the trajectory \eqref{eq:traj} records only the
\emph{even} $x$-derivatives of the flow,
\begin{equation}\label{eq:evenonly}
  \varphi^{(n)}(s)=(-1)^n\,\frac{H_{\Lambda+s}^{(2n)}(x_*)}{D}\,;
\end{equation}
hence $\varphi$ --- and every functional of it: the moments
$M_m=(-1)^mH_\Lambda^{(2m)}(x_*)$, the clock $R(\beta)$ (recovered from the
integer moments by Carlson's theorem, $M(\beta)$ being of exponential type
$0$ and bounded on $i\R$), the germ \eqref{eq:germ}, the frequency law
\eqref{eq:freq} --- depends on $H_\Lambda$ only through the even-derivative
sequence $\{H_\Lambda^{(2n)}(x_*)\}_{n\ge0}$. But real-rootedness is not an
even-derivative functional. The Laguerre--P\'olya conditions couple
consecutive derivatives, and the left--right asymmetry of the zero field
about $x_*$ --- the very two-sidedness that \eqref{eq:LPform} pinned as the
cause of the oscillation --- lives in the \emph{odd} derivatives
$\{H_\Lambda^{(2n+1)}(x_*)\}$, which $\varphi$ discards: two functions
agreeing in their even derivatives at $x_*$, one Laguerre--P\'olya and one
not, yield the same clock. No object of the program can therefore decide
$\Lambda\le0$. This is the one phenomenon behind every shot --- the cone
closing, the germ being forced, the self-reference washing out, the global
object reducing to $\Xi(x_*/2)$ without contradiction: each is the program
reading magnitude data that is blind to real-rootedness. What would decide
the Riemann hypothesis is an object sensitive to the odd, asymmetry data ---
the positions of the zeros on the two sides of $x_*$, not the value of the
flow through it.
\end{remark}

\subsection{The odd data: the asymmetry, and where the dynamics takes over}
\label{subsec:odd}
\Cref{rem:blind} names what is missing; there is a canonical object that
supplies it. Partner the value-trajectory \eqref{eq:traj} with the
\emph{slope trajectory}
\begin{equation}\label{eq:slopetraj}
  \psi(s):=H_{\Lambda+s}'(x_*),\qquad
  \psi^{(n)}(s)=(-1)^nH_{\Lambda+s}^{(2n+1)}(x_*),
\end{equation}
which records exactly the odd derivatives $\varphi$ discards, with $\psi(0)=0$
and $\psi'(0)=-H_\Lambda'''(x_*)$. Its leading datum is not another blind
oscillation: it is the global asymmetry of the zero field. Following the
midpoint $m(s)$ of the pair that the collision unfolds --- the cubic
correction to $(x-x_*)^2=2s$ shifts both roots equally --- and matching
against the de Bruijn--Newman repulsion $\dot x_j=\sum_{k\ne j}2/(x_j-x_k)$
gives
\begin{equation}\label{eq:drift}
  \dot m(0)=\frac{2}{3}\,\frac{H_\Lambda'''(x_*)}{H_\Lambda''(x_*)}
  =2\sum_{k}\frac{1}{x_*-x_k}=:2S,
\end{equation}
the Hadamard regular part of $H_\Lambda'/H_\Lambda$ at $x_*$, the net
inverse-distance to the other zeros. So the odd datum $H_\Lambda'''(x_*)$
\emph{is} the lopsidedness of the zero field about the collision --- the net
force the rest of the spectrum exerts on the colliding pair. On the zeta
zeros this $S$ is concrete and sign-changing: $+0.30$ at $x=30$ (the ordinate
$2\gamma_1=28.3$ pulling from just below), $-0.25$ at $x=35$ (pulled up by
$2\gamma_2=42.0$), $+0.69$ at $x=43$ (the near ordinate below dominating) ---
the asymmetry the even program cannot see, read straight off the neighbouring
ordinates.

This is where the deformation program ends and the dynamics begins. The
quantity $S$ --- the repulsion, the asymmetry --- is of exactly the type
whose collective behavior the zero-flow analysis of Rodgers and
Tao~\cite{RodgersTao} controls to prove $\Lambda\ge0$: their bound is a
statement about these forces, not about any value-functional. The two halves
meet at \eqref{eq:drift}: the even-data deformation program maps the
obstruction to perfect precision and is blind to it; the odd-data dynamics
touches the obstruction and has so far pinned it only from below,
$\Lambda\ge0$. The Riemann hypothesis is the matching upper control --- that
the asymmetry $S$ at a threshold collision cannot take the form a genuine
$\Lambda>0$ landing would require. That is no longer a deformation question.
The program, having carried the even data to its wall, hands the problem
across \eqref{eq:drift} to the dynamics of $S$, which is where $\Lambda\le0$
must now be won.

\subsection*{Conclusion}
Three routes to the Riemann hypothesis through generator deformation are now
resolved, negatively and for complementary reasons. The logarithmic tangent
of~\cite{Fsymbol} fails because its invariant $B_*$ is forced strictly
positive --- too definite to vanish. The even pseudo-process hierarchy of
this paper fails because its invariants $M_m$ have no fixed sign --- too
oscillatory to cone; the two are one phenomenon, the endpoint dichotomy, the
contour segment forcing a sign on the singular directions and falling silent
on the smooth even ones, with the even integers sitting at the zeros of the
imaginary-axis density that runs the program, and the level-$p$ Clark--Ocone
calculus~\cite{Fpseudo} confirming that the orders which would carry a cone
have no positivity to carry it with. And the quartic flow, the one even
direction that keeps genuine positivity, fails \emph{for} that positivity:
\cref{subsec:quartic} shows its Gaussian average over signed times fills the
dips of the de Bruijn--Newman flow and lifts conjugate pairs off the axis, so
the average manufactures nonreal zeros rather than removing them
($\Lambda(4)=\infty$ conjecturally, \cref{prob:quartic}). Too definite, too
oscillatory, too positive.

One structural cause underlies all three. By \cref{rem:blind} every object the
program builds at a collision --- moment, clock, germ, trajectory --- is a
functional of the even $x$-derivatives $\{H_\Lambda^{(2n)}(x_*)\}$ alone, and
real-rootedness is not such a functional: it is decided by the odd, asymmetry
data, the positions of the zeros on the two sides of $x_*$, which the even
program discards. The program is complete as a map of the obstruction and
constitutively unable to touch it.

What the map delivers, in the end, is the handoff. The odd datum the program
discards is not inert: by \eqref{eq:drift} it is the asymmetry
$S=\sum_k1/(x_*-x_k)$, the net repulsion the rest of the zero field exerts on
the colliding pair --- exactly the dynamical quantity whose statistics
Rodgers and Tao~\cite{RodgersTao} control to obtain $\Lambda\ge0$. Generator
deformation, carried to its wall, hands the Riemann hypothesis across the
midpoint-drift identity \eqref{eq:drift}, from the value data it exhausts to
the asymmetry $S$, on whose ledger $\Lambda\ge0$ already sits and
$\Lambda\le0$ remains to be won. That is the honest endpoint: not a proof, but
a complete account of how far deformation reaches, why it stops exactly there,
and the one identity across which the problem passes to the dynamics.

\section*{Acknowledgements}
AI tools were used as interactive assistants for drafting, reorganizing, and
generating mathematical arguments and computations, under the author's
direction. The author is responsible for all mathematical content and any
errors. Simulation code and raw output for the computational claims
accompany the source.

%==============================================================
\begin{thebibliography}{99}

\bibitem{Fsymbol} R.~Fontes, \emph{Symbol deformations of the de
Bruijn--Newman flow: infrared catastrophe, collision sensitivity, and
even-exponent rigidity}, preprint, 2026.

\bibitem{Fpseudo} R.~Fontes, \emph{A level-$p$ Clark--Ocone calculus for
higher-order heat pseudo-processes}, preprint, 2026.

\bibitem{deBruijn} N.~G.~de Bruijn, \emph{The roots of trigonometric
integrals}, Duke Math.\ J.\ \textbf{17} (1950), 197--226.

\bibitem{Newman} C.~M.~Newman, \emph{Fourier transforms with only real
zeros}, Proc.\ Amer.\ Math.\ Soc.\ \textbf{61} (1976), 245--251.

\bibitem{CSV} G.~Csordas, W.~Smith and R.~S.~Varga, \emph{Lehmer pairs of
zeros, the de Bruijn--Newman constant $\Lambda$, and the Riemann
hypothesis}, Constr.\ Approx.\ \textbf{10} (1994), 107--129.

\bibitem{RodgersTao} B.~Rodgers and T.~Tao, \emph{The de Bruijn--Newman
constant is non-negative}, Forum Math.\ Pi \textbf{8} (2020), e6.

\bibitem{Polymath15} D.~H.~J.~Polymath, \emph{Effective approximation of
heat flow evolution of the Riemann $\xi$-function, and a new upper bound for
the de Bruijn--Newman constant}, Res.\ Math.\ Sci.\ \textbf{6} (2019),
no.~31.

\bibitem{Polya27} G.~P\'olya, \emph{\"Uber trigonometrische Integrale mit
nur reellen Nullstellen}, J.\ Reine Angew.\ Math.\ \textbf{158} (1927),
6--18.

\bibitem{Titchmarsh} E.~C.~Titchmarsh, \emph{The Theory of the Riemann
Zeta-function}, 2nd ed., revised by D.~R.~Heath-Brown, Oxford University
Press, 1986.

\bibitem{Erdelyi} A.~Erd\'elyi, \emph{Asymptotic Expansions}, Dover, 1956.

\bibitem{Krylov} V.~Yu.~Krylov, \emph{Some properties of the distribution
corresponding to the equation $\partial u/\partial t=(-1)^{q+1}
\partial^{2q}u/\partial x^{2q}$}, Soviet Math.\ Dokl.\ \textbf{1} (1960),
760--763.

\bibitem{Hochberg} K.~J.~Hochberg, \emph{A signed measure on path space
related to Wiener measure}, Ann.\ Probab.\ \textbf{6} (1978), 433--458.

\bibitem{Funaki} T.~Funaki, \emph{Probabilistic construction of the solution
of some higher order parabolic differential equation}, Proc.\ Japan Acad.\
Ser.\ A Math.\ Sci.\ \textbf{55} (1979), 176--179.

\bibitem{GORZ} M.~Griffin, K.~Ono, L.~Rolen and D.~Zagier, \emph{Jensen
polynomials for the Riemann zeta function and other sequences}, Proc.\
Natl.\ Acad.\ Sci.\ USA \textbf{116} (2019), 11103--11110.

\end{thebibliography}

\end{document}
